\chapter{Introduction}
\label{chap:intro}

\section{Motivation}

Semiconductor materials are basics of modern world technology. Their usage is wide and spreads for most of the aspects in human live. Wherever to look, one can find semiconductors put to use. Cell phones, computers, laptops, servers, photovoltaic panels, microcontrollers, etc. Most commonly, they are electric circuits and electronic devices, both overtaken by silicon technology. Despite the fact that there are many more suitable, in case of physical and chemical properties, materials for electronic applications, none has dethroned silicon. Although some of them, e.g. graphene, have been hailed as the successor to silicon, until now, no material has emerged as a substitute and for a good reason. Silicon has a very well-developed technology process. Czochralski's method enabled mass production of silicon wafers with unparalleled purity and low defects rate. Moreover, at low cost. Such attributes still keeps silicon as a base material for electronic devices.

The situation differs a lot in the case of optoelectronics which is based on photons as carriers of the information, rather than electrons. Such a solution requires different physical properties of the material from those in the case of conventional electronics. The materials from which the devices are to be built must be able to emit and detect light with high efficiency. Because of that silicon, due to the indirect fundamental band gap, has no use in this area of science. Obviously, there are suitable materials to do so, but one must remember that optoelectronics is not standalone. It requires a connection to conventional electronics. That is where the problem lies. Integrating optoelectronics with electronics is not easy. It involves usage of other type of devices that will convert photonic information to electronic one. These elements are blocky and do not allow for onboard integration directly printed into the electric circuit. Moreover, their electronic properties must be very strict. Not only is a direct fundamental band gap required but also a high carrier mobility due to the fact that photons are a lot faster than the electrons. This implies that the information conversion from photonic to electronic one must be fast enough to not bottleneck the speed of photonic information flow.

The motivation for this work is to expand the capabilities and versatility of silicon technology, by finding new materials. Ones that will be integral with the silicon technology process and meet the optoelectronic devices requirements. The range of materials suitable for this task is wide. If fact to wide to focus on all of them. That is why a choice of certain group of materials must be made. This work focused on group IV elements. Main argument is easy integration with silicon technology. Carbon, germanium, or tin are isovalent with silicon, which means that they possess the same number of valence electrons. This allows to adapt technological process used with silicon to them as well. At first glance, it might seem that those materials have already been thoroughly investigated and their physical properties are well known. Most of the papers that can be found are quite old. Since then technology has improved significantly, allowing us to produce materials of lower dimensionality, higher precision of admixtures or their placement in lattice. This opens new possibilities and a chance to obtain the material composed of group IV elements that will meet optoelectronics requirements. As we know, basic group IV materials do not meet the given expectations. To achieve desirable properties, electronic structure engineering must be implemented.
\ornament
\section{Band structure engineering}

There are mainly 4 ways to make modifications in electronic structure of a material. Those are:
\begin{enumerate}
	\item strains 
	\item alloying
	\item structural factor
	\item quantum confinement
\end{enumerate}

Strains are relatively easy to implement, e.g. by chemical pressure from lattice mismatch during material growth on different bases. Depending on the material used and the type of stains, it has different effects on the electronic structure of a material. They can be isotropic, planar, or axial. The isotropic ones do not change the symmetry of the crystal. Instead, it extends or shortens the length of lattice vectors, but does not change angles between them. Planar or axial strains usually break the crystal symmetry, resulting not only in band gap modifications but may also result in band inversions or band splittings at high symmetry points in \BZ.

Alloying is another way to obtain different properties, usually by mixing a few materials or doping one with another. Often, the outcome material's parameters are a mixture of properties of mixing components. Although rarely the feature we investigate behaves linearly with doping rate. In many cases the curve is rather of a parabolic type, which also may lead to new properties of the alloy. With this method of modification, some symmetry breaks as well, due to atom substitutions or additions in the lattice. Moreover in case of low doping rate and break of periodicity of crystal, band tend to blur and new band may emerge in process.

Structural changes are more difficult to obtain. This method is basically a synthesis of the same elements but in a different space group system. Obviously, this must lead to new properties of the outcome material due to e.g. change in symmetry of the crystal. This method drastically modifies the parameters of a material, but has significant limitation. The system considered must be thermodynamically stable in different configurations, which is not always possible.

Quantum confinement, the last method presented, is based on revealing the quantum nature of a subject. That can be achieved in many ways, usually by shrinking the dimension number of a material. The surface or thin film will provide properties different from those of the bulk. The same situation occurs in the case of nanowires or dots. Any kind of discontinuity small enough to show off quantum effects may lead to enormous changes in electronic structure, e.g. bands can fold straightening the fundamental gap.

All methods presented above were implemented to conduct theoretical investigation of group IV materials. Due to the complication of structures and the periodicity of the systems, the perfect tool to use to carry out the calculations appears to be \DFT\ based programs. \abinitio\ methods are well developed and expanded to such extend that nowadays there are only few condensed matter physics theoretical works where there are no \DFT\ calculations performed or some parameters from \DFT\ used.